% Diffusion capacity theory insert for paper/draft_v1.tex.
% Recommended placement: inside Section \ref{sec:framework}, immediately after the first paragraph.

\subsection{Diffusion capacity as an economic complement}
\label{subsec:diffusion_capacity_complement}

Diffusion capacity is the institutional complement that determines whether frontier capability becomes productive capability. In a standard production view, a new technology raises output only through its interaction with other inputs. For AI in manufacturing, those inputs include sensors, robotics, machine tools, industrial data infrastructure, enterprise software, process knowledge, worker training, managerial routines, and the ability to reorganize production around prediction, inspection, control, and scheduling. A frontier model can expand the set of feasible tasks, but the realized economic effect depends on whether firms and cities possess the complementary assets needed to deploy those tasks in production.

The distinction between AI capability and AI diffusion capacity is especially important because AI is a general-purpose technology with heterogeneous local complements. A model or algorithm can be copied across places, yet the ability to use it productively depends on equipment, data, engineering talent, production culture, local suppliers, sectoral structure, and regulatory or procurement channels. The same technical capability can therefore produce different adoption outcomes across cities. The paper's empirical design treats this variation as the object of study rather than as a nuisance in the estimation problem.

The diffusion-state framework can be interpreted as a coordination technology for complementary investment. Firms may underinvest in AI-enabled production systems when standards are uncertain, when suppliers are fragmented, when the value of adoption depends on ecosystem-wide learning, or when complementary infrastructure is missing. Local governments and national ministries can alter this environment by designating pilot zones, recognizing exemplary factories, building demonstration platforms, coordinating industrial parks, and making adoption visible to firms, financiers, and officials. These channels do not eliminate selection into stronger cities. They explain why selection itself is part of the architecture of diffusion.

Hubs emerge naturally in this framework because complementary assets are spatially clustered. Cities with dense manufacturing networks, universities, engineering labor, leading firms, industrial-service providers, and experienced local governments have lower coordination costs for adopting complex production technologies. Their firms can draw on thicker supplier markets, better implementation partners, richer data environments, and stronger local policy support. A hub-centered diffusion architecture is therefore not an anomaly. It is the expected spatial form of a general-purpose technology whose productive use requires many co-located complements.

The economic role of the state in this framework is to make adoption environments more coordinated and more legible. Pilot-zone designation identifies places where national strategy, local capacity, and industrial opportunity are expected to reinforce one another. Smart-factory recognition makes selected adoption events visible and comparable across cities. Standards, demonstrations, and administrative classifications reduce uncertainty about what counts as advanced production practice. These mechanisms make the diffusion state a measurement object because they leave institutional traces of the transition from capability to use.

This framework also explains why the paper avoids a narrow treatment-effect interpretation. A policy that targets high-capacity cities, coordinates local complements, and makes adoption legible will generate patterns that combine selection and implementation. Treating that pattern as a clean treatment effect would miss the economic object. The relevant question is how a state organizes the geography of absorption and whether recognized adoption appears where the complements for industrial AI are concentrated. The hub-attenuation design, city typologies, sectoral exposure analysis, and export-relevance tables are organized around this question.
